\section{The path-sum bound}\label{sec:master}

\motiv{\vet{The commutator $\comm{A(t)}b$ of an evolved observable with a probe $b$ on a region
$Y$ can only be nonzero through the terms of the generator that cross the boundary of $Y$ (a
term inside $Y$ cannot act on an observable supported outside it): with the crossing terms
switched off, the observable's support never grows into $Y$, and causality kills the commutator.
Duhamel's identity isolates exactly the contribution of the switched-off terms, as a time
integral; hypothesis (G1) then replaces the probe $b$ by a probe on the patch $S_i$ of a single
term, one step closer to the observable, and hypothesis (G3) guarantees that the free evolution in
between cannot amplify anything.

Iterating, the commutator is bounded by a sum over paths of
terms leading from $Y$ to $X$, each weighted by the product of its strengths and by the factor $t^k/k!$ that the
nested time integrals produce. \vetnew{This chapter proves that bound,
Theorem~\ref{thm:master}, in five steps: the two analytic lemmas (Duhamel's identity and the
invariance of closed subspaces under the exponential), the support of the restricted dynamics
(a term inside the probe's region annihilates the observables outside it), the one-step
inequality, the iteration with an explicit remainder, and the limit.} Figure~\ref{fig:recursion} shows the mechanism.}}

\needspace{10\baselineskip}
\subsection{Two lemmas on the exponential}\label{subsec:exp}

\begin{lemma}[Duhamel's identity]\label{lem:duhamel}
Let $\mathbb B$ be a complete normed unital real algebra, $R,Q\in\mathbb B$ and $t\in\R$. Then
\begin{equation}\label{eq:duhamel}
  e^{t(R+Q)}=e^{tR}+\int_0^te^{(t-s)R}\,Q\,e^{s(R+Q)}\,ds,
\end{equation}
the integrand being continuous in $s$.
\end{lemma}

\begin{proof}
Let $f(s):=e^{(t-s)R}e^{s(R+Q)}$. For every \vetnew{$G\in\mathbb B$ the map $u\mapsto e^{uG}$ is
differentiable with derivative $Ge^{uG}=e^{uG}G$}, so by the chain rule and the product rule
\[
  f'(s)=-e^{(t-s)R}R\,e^{s(R+Q)}+e^{(t-s)R}(R+Q)\,e^{s(R+Q)}=e^{(t-s)R}\,Q\,e^{s(R+Q)},
\]
a continuous function of $s$. The fundamental theorem of calculus gives
$f(t)-f(0)=\int_0^tf'(s)\,ds$, and $f(t)=e^{t(R+Q)}$, $f(0)=e^{tR}$.
\end{proof}

\begin{lemma}[invariance]\label{lem:invariance}
Let $E$ be a \vetnew{complex Banach space, $M\subseteq E$ a closed complex subspace and $R$ a
bounded complex-linear operator} with $R(M)\subseteq M$. Then $e^{tR}(M)\subseteq M$ for every $t\in\R$.
\end{lemma}

\begin{proof}
For $m\in M$ the partial sums $\sum_{n\le N}t^nR^nm/n!$ lie in $M$, since $R^nm\in M$ by
induction, and they converge in norm to $e^{tR}m$. $M$ is closed.
\end{proof}

\subsection{The restricted dynamics and the first term}\label{subsec:restricted}

\vet{Fix a causal net and a family of local generators. For a region $Y$, write
$L^{\partial Y}:=L_{I^{\partial Y}}$ for the terms crossing $Y$ and
$L':=L_{I\setminus I^{\partial Y}}$ for the rest, the terms disjoint from $Y$ or inside $Y$, so
that $L=L^{\partial Y}+L'$ (Proposition~\ref{prop:subfamilies}).}

\motiv{\vet{\vetnew{The terms inside $Y$ belong with the free part}: an
observable supported outside $Y$ commutes with everything on a
patch inside $Y$, so its commutator seminorm on that patch vanishes, and (G1) then says the term
does nothing to it. Such a term cannot carry a signal out of $Y$. Only the terms that straddle
the boundary of $Y$ can, and they are the perturbation.}}

\begin{proposition}[a term inside the probe's region annihilates the complement]\label{prop:annihilate}
Let the net be causal, $S_i\subseteq Y$ and $c\in\cA(Y^c)$. Then $L_ic=0$.
\end{proposition}

\begin{proof}
Every $b\in\cA(S_i)\subseteq\cA(Y)$ commutes with $c$, since $Y^c$ and $Y$ are disjoint, so
$N_{S_i}(c)=0$ by Proposition~\ref{prop:commnorm}(b), and (G1) gives $\norm{L_ic}\le\kappa_i\cdot0$.
\end{proof}

\begin{proposition}[support of the restricted dynamics]\label{prop:support}
\begin{enumerate}[label=(\alph*),nosep]
\item Let $J\subseteq I$ and let $W$ be a region with $S_i\subseteq W$ for every $i\in J$. Then
  $L_J\,\cA(W)\subseteq\cA(W)$ and $e^{tL_J}\,\cA(W)\subseteq\cA(W)$ for every $t\in\R$.
\item Let the net be causal. Then $L'\,\cA(Y^c)\subseteq\cA(Y^c)$ and
  $e^{tL'}\,\cA(Y^c)\subseteq\cA(Y^c)$ for every $t\in\R$.
\end{enumerate}
\end{proposition}

\begin{proof}
\begin{enumerate}[label=(\alph*),nosep]
\item For $c\in\cA(W)$ and $i\in J$, (G2) gives $L_ic\in\cA(W\cup S_i)=\cA(W)$, and a finite sum of
elements of $\cA(W)$ lies in it. Lemma~\ref{lem:invariance} applies to the closed subspace
$\cA(W)$.
\item A term $i\notin I^{\partial Y}$ either misses $Y$, so that $S_i\subseteq Y^c$ and
(G2) keeps $\cA(Y^c)$ as in (a), or lies inside $Y$, and then $L_ic=0$ for $c\in\cA(Y^c)$ by
Proposition~\ref{prop:annihilate}. Summing, $L'$ maps $\cA(Y^c)$ into itself, and
Lemma~\ref{lem:invariance} applies.\qedhere
\end{enumerate}
\end{proof}

\begin{proposition}[the first term vanishes]\label{prop:first}
Let the net be causal. If $X\cap Y=\emptyset$, $a\in\cA(X)$ and $b\in\cA(Y)$, then
$\comm{e^{tL'}a}{b}=0$ for every $t\in\R$.
\end{proposition}

\begin{proof}
$X\subseteq Y^c$, so $a\in\cA(Y^c)$ and $e^{tL'}a\in\cA(Y^c)$ by
Proposition~\ref{prop:support}(b); $Y^c$ and $Y$ are disjoint, so causality gives the claim.
\end{proof}

\subsection{The one-step inequality}\label{subsec:onestep}

\needspace{10\baselineskip}
\begin{proposition}[the reset step]\label{prop:reset}
For $s\le t$, $c\in\cA$ and $b\in\cA$ with $\norm b\le1$,
\[
  \bigl\lVert\comm{e^{(t-s)L'}\,L^{\partial Y}c}{b}\bigr\rVert\le2\sum_{i\in I^{\partial Y}}\kappa_i\,N_{S_i}(c).
\]
\end{proposition}

\begin{proof}
$\norm{\comm{c'}b}\le2\norm{c'}$ for every $c'\in\cA$ and $\norm b\le1$; $\norm{e^{(t-s)L'}}\le1$
by (G3) for the sub-family $I\setminus I^{\partial Y}$ and $t-s\ge0$; and
$\norm{L^{\partial Y}c}\le\sum_{i\in I^{\partial Y}}\kappa_iN_{S_i}(c)$ by
Proposition~\ref{prop:subfamilies}(c).
\end{proof}

\begin{figure}[t]
\centering
% Figure 2: the probe-reset recursion. Transcription (for cheap agents): a tree. Root node
% "C_{X,Y}(t)" (commutator of A(t) on X with a probe on Y). One Duhamel step splits it into
% "2[X meets Y]" (the trivial bound when the regions meet; for disjoint regions the free term of
% the Duhamel step is zero) plus, for each term i crossing Y (meeting Y and its complement), a
% branch labelled "2 kappa_i int_0^t ... ds" leading to the node "C_{X,S_i}(s)" (the same
% commutator with a probe on S_i). Branches whose S_i meets X stop with the trivial bound "<= 2";
% the others split again in the same way. The sequence of patches along a root-to-leaf branch is
% a first-contact path. Each branch label is written horizontally near the lower end of its
% arrow, on the outer side (left of a left branch, right of a right branch), clear of the other
% arrows; the free-term node sits higher, to the far left, with its note below it. The dashed
% arrow to the further branches leaves the root to the right and turns down to the dots.
\begin{tikzpicture}[x=1cm,y=1cm,
  node/.style={draw,rounded corners=3pt,inner sep=4pt,font=\small,fill=white},
  leaf/.style={node,draw=Blue,fill=Blue!6},
  lab/.style={font=\scriptsize,fill=white,inner sep=1pt},
  >={Stealth[length=5pt]}]
\node[node] (root) at (0,0) {$C_{X,Y}(t)$};
\node[node] (free) at (-5.4,-1.0) {$2\,\ind{X\cap Y\ne\emptyset}$};
\node[lab,below=1pt of free,align=center] {free term\\ $=0$ when $X\cap Y=\emptyset$};
\node[node] (i1) at (-1.9,-2.0) {$C_{X,S_{i_1}}(s)$};
\node[node] (i2) at (1.9,-2.0) {$C_{X,S_{i_2}}(s)$};
\node[font=\small] (dots) at (4.8,-2.0) {$\cdots$};
\draw[->] (root) -- (free);
\draw[->] (root) -- (i1) node[lab,pos=0.82,above left=1pt]{$2\kappa_{i_1}\!\int_0^t\!ds$};
\draw[->] (root) -- (i2) node[lab,pos=0.82,above right=1pt]{$2\kappa_{i_2}\!\int_0^t\!ds$};
\draw[->,dashed] (root.east) -| (dots.north);
% second level under i1: one branch stops (meets X), one continues
\node[leaf] (stop) at (-3.9,-4.0) {$\le2$ \ ($S_{j}\cap X\ne\emptyset$)};
\node[node] (cont) at (0.1,-4.0) {$C_{X,S_{j'}}(s')$};
\draw[->] (i1) -- (stop) node[lab,pos=0.82,above left=1pt]{$2\kappa_{j}\!\int_0^s\!ds'$};
\draw[->] (i1) -- (cont) node[lab,pos=0.82,above right=1pt]{$2\kappa_{j'}\!\int_0^s\!ds'$};
\draw[->,dashed] (cont) -- ++(0,-1.0) node[below,font=\scriptsize]{continues until a term meets $X$};
\node[font=\scriptsize,align=left,anchor=west] at (2.4,-4.0)
  {a root-to-leaf branch $Y\to S_{i_1}\to S_j$\\ is a first-contact path; its weight\\
   $\prod 2\kappa$, the leaf bound $2$ and the\\ nested integrals give $2\,(2t)^k/k!\,\prod\kappa$};
\end{tikzpicture}

\caption{\vet{The probe-reset recursion. The commutator of the evolved observable $A(t)$ on $X$
with a probe on $Y$ is bounded, after one Duhamel step, by the terms $S_i$ crossing $Y$, each
contributing $2\kappa_i$ times a time integral of the same commutator with a probe on $S_i$. The
recursion stops at the first term meeting $X$, where the trivial bound $2$ is used; the chain of
patches visited is a path of Definition~\ref{def:path}.}}
\label{fig:recursion}
\end{figure}

\vet{The reset step is where the probe moves: on the right-hand side the commutator with the
original probe $b$ has become the seminorm $N_{S_i}(c)$, a supremum of commutators with probes
on the single patch $S_i$. The next statement packages one application of Duhamel's identity
with this step. It takes an arbitrary continuous majorant of the reset seminorms, because that
is how the induction will consume it. The theorem needs the regions disjoint, so a probe on a
patch that meets $X$ cannot be reset again; there the trivial bound $2$ of
Proposition~\ref{prop:commnorm}(a) is used instead. This is why paths end at first contact.}

\begin{theorem}[one Duhamel step]\label{thm:onestep}
Let the net be causal, $X\cap Y=\emptyset$, $a\in\cA(X)$, $b\in\cA(Y)$ with $\norm b\le1$,
$t\ge0$, and let $\beta_i:\R\to\R$ ($i\in I^{\partial Y}$) be continuous functions with
$N_{S_i}(A(s))\le\beta_i(s)$ for all $s\in[0,t]$. Then
\begin{equation}\label{eq:onestep}
  \norm{\comm{A(t)}b}\le2\sum_{i\in I^{\partial Y}}\kappa_i\int_0^t\beta_i(s)\,ds .
\end{equation}
\end{theorem}

\begin{proof}
Lemma~\ref{lem:duhamel} with $R=L'$ and $Q=L^{\partial Y}$, applied to $a$, gives
$A(t)=e^{tL'}a+\int_0^te^{(t-s)L'}L^{\partial Y}A(s)\,ds$. Evaluation at $a$ and the commutator
with $b$ are bounded linear maps and pass through the integral:
\[
  \comm{A(t)}b=\comm{e^{tL'}a}{b}+\int_0^t\comm{e^{(t-s)L'}L^{\partial Y}A(s)}{b}\,ds .
\]

The first term is $0$ by Proposition~\ref{prop:first}. The norm of the integral is at most the
integral of the norm, which Proposition~\ref{prop:reset} with $c=A(s)$ bounds pointwise by
$2\sum_{i\in I^{\partial Y}}\kappa_iN_{S_i}(A(s))\le2\sum_{i\in I^{\partial Y}}\kappa_i\beta_i(s)$,
a continuous function.
\end{proof}

\subsection{The iteration and the series}\label{subsec:iteration}

\vet{For $k\ge0$ and $t\in\R$ put $\gamma_k(t):=(2t)^{k+1}/(k+1)!$, and for regions $X,Y$,
$n\ge0$ and $t\in\R$ (a polynomial in $t$, used for $t\ge0$),
\begin{equation}\label{eq:truncbound}
  B_n(Y,X,t):=2\Bigl(\sum_{k<n}\gamma_k(t)\,P_{k+1}(Y,X)+\gamma_n(t)\,K^{n+1}\Bigr),
\end{equation}
the truncated series with its remainder. The identity $\int_0^t\gamma_k(s)\,ds=\gamma_{k+1}(t)/2$
is what makes the time-ordered integrals produce the factorials.}

\begin{theorem}[iteration]\label{thm:iteration}
Let the net be causal, $X\cap Y=\emptyset$, $a\in\cA(X)$, $b\in\cA(Y)$ with
$\norm a,\norm b\le1$, and $t\ge0$.
Then $\norm{\comm{A(t)}b}\le B_n(Y,X,t)$ for every $n\ge0$.
\end{theorem}

\begin{proof}
Induction on $n$, for all disjoint pairs $X,Y$, all $t\ge0$ and all $a,b$ of norm at most $1$
at once.

\emph{Base.} Theorem~\ref{thm:onestep} with the constant majorants $\beta_i=2$, legitimate since
$N_{S_i}(A(s))\le2\norm{A(s)}\le2$ by Propositions~\ref{prop:commnorm}(a)
and~\ref{prop:subfamilies}(b), gives $\norm{\comm{A(t)}b}\le4t\sum_{i\in I^{\partial Y}}\kappa_i\le4tK
=B_0(Y,X,t)$.

\emph{Step.} Assume the bound for $n$. Apply Theorem~\ref{thm:onestep} with the majorant
$\beta_i(s)=2$ when $S_i$ meets $X$, and $\beta_i(s)=B_n(S_i,X,s)$ when $S_i$ misses $X$; the
latter is
a majorant because, for every unit $b'\in\cA(S_i)$, the induction hypothesis for the disjoint
pair $(X,S_i)$ gives $\norm{\comm{A(s)}{b'}}\le B_n(S_i,X,s)$, and $N_{S_i}(A(s))$ is the
supremum of these. Integrating \eqref{eq:truncbound} with the identity for $\int\gamma_k$,
\[
  \int_0^tB_n(S_i,X,s)\,ds=\sum_{k<n}\gamma_{k+1}(t)\,P_{k+1}(S_i,X)+\gamma_{n+1}(t)\,K^{n+1}.
\]

Summing over $i\in I^{\partial Y}$ with the weights $2\kappa_i$: the terms meeting $X$ contribute
$2\gamma_0(t)\sum\kappa_i=2\gamma_0(t)P_1(Y,X)$; the terms missing $X$ contribute
$2\sum_{k<n}\gamma_{k+1}(t)\sum_i\kappa_iP_{k+1}(S_i,X)=2\sum_{k<n}\gamma_{k+1}(t)P_{k+2}(Y,X)$ by
the recursion of Proposition~\ref{prop:pathcount}(a), and a remainder
$2\gamma_{n+1}(t)K^{n+1}\sum_i\kappa_i\le2\gamma_{n+1}(t)K^{n+2}$, where the sum over the terms
missing $X$ is at most $K$. Together this is $B_{n+1}(Y,X,t)$: the first group is the $k=0$
term of the sum in \eqref{eq:truncbound} for $n+1$, the second group its terms $k=1,\dots,n$,
and the remainder its last term.
\end{proof}

\begin{theorem}[the path-sum bound]\label{thm:master}
Let the net be causal and $t\ge0$. For regions $X,Y$, $a\in\cA(X)$ and $b\in\cA(Y)$,
\begin{equation}\label{eq:master}
  \norm{\comm{A(t)}b}\le2\norm a\norm b\Bigl(\ind{X\cap Y\ne\emptyset}
  +\sum_{k\ge1}\frac{(2t)^k}{k!}\,P_k(Y,X)\Bigr),
\end{equation}
where $\ind{P}$ is the indicator of a condition $P$, $1$ if it holds and $0$ otherwise; the series converges, with
sum at most $e^{2tK}-1$. Equivalently, for
$C_{X,Y}(t):=\sup\{\norm{\comm{e^{tL}a}b}:a\in\cA(X),b\in\cA(Y),\norm a,\norm b\le1\}$,
\[
  C_{X,Y}(t)\le2\,\ind{X\cap Y\ne\emptyset}+2\sum_{k\ge1}\frac{(2t)^k}{k!}\,P_k(Y,X).
\]
\end{theorem}

\begin{proof}
The $k$-th term is at most $(2tK)^k/k!$ by Proposition~\ref{prop:pathcount}(b), so the series
converges. If $X$ meets $Y$ the bound $2\norm a\norm b$ holds since $\norm{A(t)}\le\norm a$. If
$X\cap Y=\emptyset$ and $a,b$ are unit, Theorem~\ref{thm:iteration} gives, for every $n$, the
truncated series plus $2(2tK)^{n+1}/(n+1)!$, which tends to $0$; the truncated series is at most
the full one. General $a,b$ follow by scaling, and the supremum form by taking the supremum.
\end{proof}

\begin{example}[two qubits, continued]\label{ex:twosites-bound}
\vet{For $H=\sigma^z_0\sigma^z_1$, $a=\sigma^x_0$ and $b=\sigma^x_1$ (Example~\ref{ex:twosites}),
$P_1=1$ and $P_k=0$ for $k\ge2$, so Theorem~\ref{thm:master} gives
$\norm{\comm{A(t)}b}\le2\cdot2t=4t$. The exact value $2|\sin2t|$ lies below $4t$ and matches
it to first order: the constant $2\kappa$ per step is not loose at short times.}
\end{example}

\begin{example}[a three-qubit chain]\label{ex:chain}
\vet{Three qubits on $\Lambda=\{0,1,2\}$ with the Hamiltonian
\[
  H=\sigma^z_0\sigma^z_1+\tfrac12\,\sigma^x_1\sigma^x_2 ,
\]
that is, two terms: $\Phi_1=\sigma^z_0\sigma^z_1$ on $S_1=\{0,1\}$ with $\kappa_1=1$, and
$\Phi_2=\tfrac12\sigma^x_1\sigma^x_2$ on $S_2=\{1,2\}$ with $\kappa_2=\tfrac12$. The observable
is $a=\sigma^x_0$ on $X=\{0\}$, the probe $b=\sigma^z_2$ on $Y=\{2\}$.

\emph{The paths.} Only $S_2$ crosses $Y$, and it misses $X$; after it only $S_1$ can follow,
since $S_1$ crosses $S_2$ while $S_2$ does not cross itself, and $S_1$ meets $X$. So the only
path is $(S_2,S_1)$: $P_2=\kappa_1\kappa_2=\tfrac12$, every other $P_k$ is $0$, and
Theorem~\ref{thm:master} gives
\[
  \norm{\comm{A(t)}{\sigma^z_2}}\le2\cdot\frac{(2t)^2}{2!}\cdot\frac12=2t^2\qquad\vetnew{\text{for every }t\ge0}.
\]

\vetnew{\emph{The exact value.} The three Pauli strings $\mathbf e_0=\sigma^x_0$,
$\mathbf e_1=\sigma^y_0\sigma^z_1$, $\mathbf e_2=\sigma^y_0\sigma^y_1\sigma^x_2$ satisfy
$\comm H{\mathbf e_0}=2\ii\mathbf e_1$, $\comm H{\mathbf e_1}=-2\ii\mathbf e_0-\ii\mathbf e_2$,
$\comm H{\mathbf e_2}=\ii\mathbf e_1$, so
$A(t)=f_0(t)\mathbf e_0+f_1(t)\mathbf e_1+f_2(t)\mathbf e_2$ with $f_0'=2f_1$,
$f_1'=-2f_0-f_2$, $f_2'=f_1$ and $(f_0,f_1,f_2)(0)=(1,0,0)$; hence
$f_2(t)=\tfrac25(\cos\sqrt5\,t-1)$. Only $\mathbf e_2$ fails to commute with the probe,
$\comm{\mathbf e_2}{\sigma^z_2}=-2\ii\,\sigma^y_0\sigma^y_1\sigma^y_2$, so}
\[
  \norm{\comm{A(t)}{\sigma^z_2}}=\vetnew{2|f_2(t)|}=\tfrac45\bigl(1-\cos\sqrt5\,t\bigr)
  =2t^2-\tfrac56t^4+O(t^6).
\]
The ratio to the bound is $1-\tfrac5{12}t^2+\dots$, equal to $0.99996$ at $t=0.01$ and
decreasing; the bound is never exceeded. On this chain the path-sum bound is sharp to leading
order at short times, while at long times every commutator is at most $2$ and the bound grows.}
\end{example}

\subsection*{Route against the source}
\vet{The sources differentiate the commutator in time and move the derivative onto the probe
with the Jacobi identity. Writing $f(t)=\comm{\tau_t(A)}B$ and
$\frac{d}{dt}\tau_t(A)=\ii\comm{\tau_t(H_X)}{\tau_t(A)}$ with $H_X$ the terms meeting the
observable's region, \cite{nsy2019} (eqs.~(3.40)--(3.48)) and \cite{nachtergaele-sims} obtain
$f'=\ii\comm{\tau_t(H_X)}{f}-\ii\comm{\tau_t(A)}{\comm{\tau_t(H_X)}B}$; the first term generates a
norm-preserving flow and drops out of the norm estimate, the second resets the observable to a
single term $\Phi(Z)$ with $Z\cap X\ne\emptyset$, and iterating gives a series over chains of
overlapping supports. The review \cite{chen-lucas-yin} (eqs.~(3.26)--(3.29)) does the same in
graph language.

The route above differs in three respects. It integrates (Duhamel) instead of
differentiating, so that no Jacobi identity and no derivation property are needed, and the
hypotheses reduce to (G1)--(G3); the open-system bound of Chapter~\ref{sec:approx} is therefore
the same theorem. It resets the probe rather than the observable, which is a matter of taste.
And it stops at first contact, which the sources do not. The path class is otherwise theirs.

The chains of \cite{nsy2019}, eqs.~(3.41)--(3.42), are sequences $Z_1,\dots,Z_n$ with $Z_1$
meeting the observable's region and its complement, each $Z_{j+1}$ meeting $Z_j$ and its
complement, and $Z_n$ meeting the probe's region; they run from the observable to the probe
because the source resets the observable. For an automorphism group
$\norm{\comm{\tau_t(A)}B}=\norm{\comm{A}{\tau_{-t}(B)}}$ and the source's bound depends on $|t|$,
so its theorem applied to the pair $(B,A)$ runs the chains from $Y$: $Z_1$ crosses $Y$ and each
$Z_{j+1}$ crosses $Z_j$, which is Definition~\ref{def:path} without first contact. With one
term per support and $\kappa_Z=\norm{\Phi(Z)}$, the weights $2\norm{\Phi(Z_j)}$ and the
nested integrals $t^n/n!$ are the same, so in the Hamiltonian case the series \eqref{eq:master}
is at most the chain series of \cite{nsy2019} term by term.

The reason the sources may restrict
to crossing sets is the one of Proposition~\ref{prop:annihilate}: a term inside the current
region cannot move an observable supported outside it. The review's recursion (3.29b) drops the
repeated edge without giving this reason; Proposition~\ref{prop:annihilate} supplies it.

The review's Theorem~3.4 sums over sequences from a set meeting $Y$ to a set meeting $X$ with
consecutive overlap and distinct consecutive terms, a class containing ours since a set crossing
another differs from it, so \eqref{eq:master} is at most that series as well; the review
attributes the theorem to the work it numbers [34], and the stronger irreducible-path bound of
\cite{chen-lucas2021}, Theorem~3, is not claimed here.}

\subsection*{Notes and references}
\vet{Theorem~\ref{thm:master} is the generator-level form of the Lieb-Robinson path series,
Theorem~3.4 of \cite{chen-lucas-yin} and eqs.~(3.41)--(3.42) of \cite{nsy2019}, with the
differences listed above; its constants, $2\kappa_i$ per step and $t^k/k!$, are theirs.
Duhamel's identity and the invariance lemma are standard; both are proved here for a general
complete normed unital real algebra and a general Banach space.}

\paragraph{Checks.} {\small The numbers of Example~\ref{ex:chain} are computed by
\texttt{checks/m1a\_crossing\_check.py} and \texttt{checks/m1a\_chain\_expansion\_check.py},
the latter also for the closed form of the commutator norm.}
